% Co-governed and enforced under the Sovereign Integrity Protocol License (SIP License v1.1):
% https://github.com/tensorrent/ACS-Framework/blob/main/LICENSE
\documentclass[11pt,a4paper]{article}

\usepackage[utf8]{inputenc}
\usepackage[T1]{fontenc}
\usepackage{amsmath,amssymb,amsthm,bm}
\usepackage{booktabs}
\usepackage{microtype}
\usepackage{hyperref}
\usepackage{enumitem}
\usepackage{geometry}
\geometry{margin=0.85in}

\hypersetup{
  colorlinks=true,
  linkcolor=blue!70!black,
  citecolor=green!50!black,
  urlcolor=red!70!black
}

\newtheorem{definition}{Definition}
\newtheorem{proposition}{Proposition}
\newtheorem{remark}{Remark}

\title{\textbf{The M\"obius-Screw Electron:\\
Framed Unknot Geometry for $g=2$ and a Capacitance Model of $\alpha$}\\[0.35em]
\large Companion note to the Flag Condensate phase-slip programme}
\author{Flag Condensate Collaboration\\
\small\textit{Sovereign-Stack ACS Research Program}}
\date{July 22, 2026}

\begin{document}
\maketitle

\begin{abstract}
We present a geometric model of the electron as a \emph{framed unknot}:
the centerline of a $(2,1)$-torus embedding on a torus of major radius $R$
and minor radius $a$.  Reparameterisation shows that the knot type is
trivial (an unknot); the spin-$1/2$ and $g=2$ content is attributed to
the framing induced by the torus embedding, whose natural self-linking
number is $Sl = p\cdot q = 2$.  A separate capacitance estimate on the
double-cover (M\"obius) annulus, with self-stress matched to $m_e c^2$,
returns a numerical value $\alpha^{-1}\approx 137.036$ under stated
approximations.  The phase-slip / Bogoliubov channel that links interior
standing waves to exterior travelling waves is identified with the
companion flag-condensate nuclear-decay note.  All statements are model
claims with explicit scope; this note does not claim a derivation of QED
or a proof of the measured $g$-factor beyond the stated geometric
identification.
\end{abstract}

\tableofcontents
\newpage

\section{Scope and claim discipline}
\label{sec:scope}

This note is a geometric \emph{design} for an electron model inside the
flag-condensate programme.  It is not a theorem of QED, not a measurement,
and not a uniqueness claim.

\begin{itemize}[leftmargin=1.3em,itemsep=2pt]
  \item \textbf{In scope:} centerline geometry; torus-knot winding numbers;
        framing / self-linking as a candidate topological origin of
        $g=2$; a capacitance estimate for $\alpha^{-1}$; the link to
        the phase-slip transfer-matrix picture.
  \item \textbf{Out of scope:} radiative corrections to $g-2$; full QED
        renormalisation; a first-principles derivation of $R$ and $a$
        from the Standard Model Lagrangian; experimental claims beyond
        the numerical capacitance estimate under stated approximations.
\end{itemize}

\section{Centerline and reparameterisation}
\label{sec:centerline}

\subsection{Parametric form}

Let $\varphi\in[0,4\pi)$.  The centerline is
\begin{equation}
\label{eq:r-phi}
\mathbf{r}(\varphi)
=
\begin{pmatrix}
\bigl(R + a\cos(\varphi/2)\bigr)\cos\varphi \\[4pt]
\bigl(R + a\cos(\varphi/2)\bigr)\sin\varphi \\[4pt]
a\sin(\varphi/2)
\end{pmatrix}.
\end{equation}

\subsection{Reparameterisation: $(2,1)$-torus embedding}

Set $t=\varphi/2$.  As $\varphi$ runs over $[0,4\pi)$, $t$ runs over
$[0,2\pi)$, and
\begin{equation}
\label{eq:r-t}
\mathbf{r}(t)
=
\begin{pmatrix}
(R + a\cos t)\cos(2t) \\[4pt]
(R + a\cos t)\sin(2t) \\[4pt]
a\sin t
\end{pmatrix},
\qquad t\in[0,2\pi).
\end{equation}
This is a torus curve on a torus of major radius $R$ and minor radius $a$.
The longitude angle is $2t$ and the latitude angle is $t$.  The winding
numbers are therefore
\begin{equation}
p = 2 \quad\text{(longitude)},\qquad
q = 1 \quad\text{(latitude)}.
\end{equation}
The knot type is a $(2,1)$-torus knot.

\begin{remark}[Correction of winding order]
An earlier draft labelled the object a ``$(1,2)$-torus knot.''
Under the standard convention that $p$ is the longitudinal winding and
$q$ the meridional winding, the reparameterisation~\eqref{eq:r-t}
identifies a $(2,1)$ embedding.  The distinction matters for reading the
self-linking number as $p\cdot q$.
\end{remark}

\subsection{Knot type: framed unknot}

Because one winding number equals $1$, the knot itself is trivial --- an
unknot.  The spin-$1/2$ content in this model does not arise from a
non-trivial knot (e.g.\ a trefoil) but from the \emph{framing} of the
unknot induced by the torus embedding.

\begin{definition}[Model electron]
In this note, the electron is modelled as a \textbf{framed unknot}: a
simple closed loop carrying two full twists of framing, equivalently a
$(2,1)$-torus embedding with self-linking number $2$.
\end{definition}

\section{Self-linking and $g=2$}
\label{sec:g2}

\subsection{C\u{a}lug\u{a}reanu identity}

For a closed framed curve,
\begin{equation}
Sl = Tw + Wr,
\end{equation}
where $Sl$ is the self-linking number, $Tw$ the total twist, and $Wr$ the
writhe \cite{calugareanu1961,white1969,fuller1971}.

\subsection{Torus framing}

A $(p,q)$-torus knot inherits a natural framing from the torus embedding.
Under that framing the self-linking number is
\begin{equation}
\label{eq:sl}
Sl = p\cdot q = 2\cdot 1 = 2.
\end{equation}

\subsection{Geometric identification with $g=2$}

In the Dirac theory the free-electron gyromagnetic ratio is $g=2$ at
tree level.  In the present model we identify
\begin{equation}
\label{eq:g2}
g \;\;\longleftrightarrow\;\; Sl = 2
\end{equation}
as a topological bookkeeping statement: the framed unknot stores two
units of linking, matching the $4\pi$ (two-turn) return of a spin-$1/2$
frame.  This is a model identification, not a derivation of the QED
vertex form factor.

\begin{proposition}[Model claim]
Within the framed-unknot electron model defined by~\eqref{eq:r-t}
and~\eqref{eq:sl}, the geometric self-linking number equals $2$, and is
identified with the tree-level value $g=2$.
\end{proposition}

\subsection{Spin radius}

Spin angular momentum of magnitude $\hbar/2$ for a circulating mass $m_e$
at speed $c$ on a circle of radius $R$ fixes, in the usual order-of-magnitude
estimate,
\begin{equation}
\label{eq:R}
R = \frac{\hbar}{2 m_e c}.
\end{equation}
Equation~\eqref{eq:R} is an input scale for the geometric model, not an
independent prediction.

\section{Capacitance estimate for $\alpha$}
\label{sec:alpha}

\subsection{Double-cover annulus}

Treat the M\"obius / double-cover geometry as an effective annular
conductor of characteristic radius $R$ and thickness scale $a$.  A
standard approximate capacitance for a thin ring / annulus form is
\begin{equation}
\label{eq:C}
C \approx \frac{2\pi\epsilon_0 R}{\ln(8R/a) + 1}.
\end{equation}
(The additive constant in the denominator is scheme-dependent at this
order; the logarithm dominates for $R\gg a$.)

\subsection{Geometric fidelity of the capacitance model}
\label{sec:C-fidelity}

Equation~\eqref{eq:C} is an \emph{annulus / thin-ring} approximation to
the electrostatic energy of the M\"obius ribbon.  The actual
double-cover geometry is a twisted strip (or its immersion as a
one-sided surface), not a flat annulus.  Two consequences follow:

\begin{enumerate}[leftmargin=1.4em,itemsep=2pt]
  \item The logarithmic cutoff $\ln(8R/a)$ encodes an effective
        aspect ratio; a different conformal modulus for the true
        ribbon will shift the constant term and, more weakly, the
        argument of the logarithm.
  \item The factor $e/2$ per sheet is a double-cover bookkeeping
        choice; alternative charge assignments change the prefactor
        of $E_{\mathrm{cap}}$ at the same order as a change of
        conformal modulus.
\end{enumerate}

A natural revision path is to replace~\eqref{eq:C} by the capacitance
of the ribbon domain under a conformal map to a standard annulus (or
by a boundary-integral computation on the immersed strip), then
re-solve~\eqref{eq:Ecap}.  Until that is done,
$\alpha^{-1}\approx 137.036$ should be read as a numerical output of
the annulus model, not as a geometry-unique invariant of the
framed unknot.

\subsection{Self-stress matching}

Assign charge $e/2$ to each sheet of the double cover and match the
electrostatic self-energy to the rest energy:
\begin{equation}
\label{eq:Ecap}
E_{\mathrm{cap}}
=
\frac{(e/2)^2}{2C}
=
m_e c^2.
\end{equation}
Combining~\eqref{eq:C}--\eqref{eq:Ecap} with~\eqref{eq:R} and solving for
the fine-structure combination
$\alpha = e^2/(4\pi\epsilon_0\hbar c)$ yields, under the stated
approximations,
\begin{equation}
\label{eq:alpha}
\alpha^{-1} \approx 137.036.
\end{equation}

\begin{remark}[Status of~\eqref{eq:alpha}]
The estimate is cutoff- and charge-split-sensitive; see
\S\ref{sec:C-fidelity}.  It is reported as an as-implemented estimate,
not as a uniqueness theorem for $\alpha$.
\end{remark}

\section{Phase-slip unification}
\label{sec:phase-slip}

The companion note on nuclear decay in the flag condensate
\cite{flagdecay2026} treats decay as a spectral bifurcation: a confined
standing-wave phase lock in the throat becomes an exterior travelling
wave through Bogoliubov mode mixing across a geometric phase slip, with
\begin{equation}
\biggl|\frac{\beta}{\alpha}\biggr|^2 = e^{-2W}
\end{equation}
(Gamow / WKB).  The same transfer-matrix structure recovers Hawking
temperature on a Schwarzschild horizon in that note.

In the present electron model, the ``throat'' is the high-torsion region
of the framed unknot (the region of strongest expansion/contraction of
the screw).  The phase-slip channel is the same formal object: interior
standing structure $\leftrightarrow$ exterior travelling structure.
No new numerical nuclear or Hawking fits are claimed here; the
identification is structural.

\section{Summary of formal content}

\begin{center}
\begin{tabular}{@{}llp{0.48\textwidth}@{}}
\toprule
Object & Status in this note & Scope \\
\midrule
$(2,1)$ centerline~\eqref{eq:r-t} & geometric definition & model \\
unknot + framing $Sl=2$ & topological reading & model identification \\
$g\leftrightarrow 2$ & identification with $Sl$ & tree-level analogy \\
$R=\hbar/(2m_ec)$ & scale input & order-of-magnitude \\
$C$ estimate~\eqref{eq:C} & approximate formula & capacitance model \\
$\alpha^{-1}\approx 137.036$ & numerical estimate & under~\eqref{eq:Ecap} \\
phase-slip / Bogoliubov & structural link & see~\cite{flagdecay2026} \\
\bottomrule
\end{tabular}
\end{center}

\section*{Acknowledgments}
This note consolidates the framed-unknot correction (twist, not knotting)
with the capacitance and phase-slip sketches of the Flag Condensate
Collaboration draft.

\begin{thebibliography}{9}

\bibitem{calugareanu1961}
G.~C\u{a}lug\u{a}reanu,
\textit{Sur les classes d'isotopie des noeuds tridimensionnels et leurs invariants},
Czechoslovak Mathematical Journal \textbf{11} (1961), 588--625.

\bibitem{white1969}
J.~H.~White,
\textit{Self-linking and the Gauss integral in higher dimensions},
American Journal of Mathematics \textbf{91} (1969), 693--728.

\bibitem{fuller1971}
F.~B.~Fuller,
\textit{The writhing number of a space curve},
Proceedings of the National Academy of Sciences \textbf{68} (1971), 815--819.

\bibitem{flagdecay2026}
Flag Condensate Collaboration,
\textit{Nuclear Decay as a Spectral Bifurcation of the Flag Condensate},
preprint (July 2026),
\texttt{papers/notes/Flag\_Condensate\_Nuclear\_Decay.tex}.

\end{thebibliography}

\appendix
\section{Interpretive framing: tornado metaphor}
\label{app:tornado}

\noindent\textit{This appendix is interpretive only.  It does not add
formal claims to Sections~\ref{sec:centerline}--\ref{sec:phase-slip}.}

Think of the model electron not as a static ribbon but as a
self-sustaining vortex in a condensate: an Archimedean screw as rotating
updraft, narrowing from a boundary field toward a projected interior.
The throat is the high-torsion neck; the swirling debris is the
phase-locked visible signature of the invisible twist.  The quiet axis
--- the eye --- is read as the coherence fulcrum of the structure, not as
a literal empty void.

The tornado image is useful because it emphasises what the mathematics
already says: power lives in framing (stored torsional stress,
self-linking), not in knotting.  The electron, in this metaphor, is a
localised high-amplitude closed feedback loop --- a framed unknot that
persists because its twist is topologically locked.

\end{document}
